% Options for packages loaded elsewhere
\PassOptionsToPackage{unicode}{hyperref}
\PassOptionsToPackage{hyphens}{url}
\documentclass[
  12pt,
]{ctexart}
\usepackage{xcolor}
\usepackage{amsmath,amssymb}
\setcounter{secnumdepth}{-\maxdimen} % remove section numbering
\usepackage{iftex}
\ifPDFTeX
  \usepackage[T1]{fontenc}
  \usepackage[utf8]{inputenc}
  \usepackage{textcomp} % provide euro and other symbols
\else % if luatex or xetex
  \usepackage{unicode-math} % this also loads fontspec
  \defaultfontfeatures{Scale=MatchLowercase}
  \defaultfontfeatures[\rmfamily]{Ligatures=TeX,Scale=1}
\fi
\usepackage{lmodern}
\ifPDFTeX\else
  % xetex/luatex font selection
\fi
% Use upquote if available, for straight quotes in verbatim environments
\IfFileExists{upquote.sty}{\usepackage{upquote}}{}
\IfFileExists{microtype.sty}{% use microtype if available
  \usepackage[]{microtype}
  \UseMicrotypeSet[protrusion]{basicmath} % disable protrusion for tt fonts
}{}
\makeatletter
\@ifundefined{KOMAClassName}{% if non-KOMA class
  \IfFileExists{parskip.sty}{%
    \usepackage{parskip}
  }{% else
    \setlength{\parindent}{0pt}
    \setlength{\parskip}{6pt plus 2pt minus 1pt}}
}{% if KOMA class
  \KOMAoptions{parskip=half}}
\makeatother
\usepackage{longtable,booktabs,array}
\usepackage{calc} % for calculating minipage widths
% Correct order of tables after \paragraph or \subparagraph
\usepackage{etoolbox}
\makeatletter
\patchcmd\longtable{\par}{\if@noskipsec\mbox{}\fi\par}{}{}
\makeatother
% Allow footnotes in longtable head/foot
\IfFileExists{footnotehyper.sty}{\usepackage{footnotehyper}}{\usepackage{footnote}}
\makesavenoteenv{longtable}
\setlength{\emergencystretch}{3em} % prevent overfull lines
\providecommand{\tightlist}{%
  \setlength{\itemsep}{0pt}\setlength{\parskip}{0pt}}
\usepackage{bookmark}
\IfFileExists{xurl.sty}{\usepackage{xurl}}{} % add URL line breaks if available
\urlstyle{same}
\hypersetup{
  hidelinks,
  pdfcreator={LaTeX via pandoc}}

\author{SCX}
\date{2026}

\begin{document}

\section{SCX TODO --- 2026-06-26}\label{scx-todo-2026-06-26}

\begin{quote}
\textbf{本文件已被 \texttt{SCX\_TODO\_v4\_通用模块化.md} 替代}
(2026-06-26)

v4 通用模块化重构是当前主方向。本文档保留仅用于历史参考，不再作为当前
TODO。 请查看 \texttt{SCX\_TODO\_v4\_通用模块化.md} 和
\texttt{SCX\_规划\_v4\_通用模块化.md}。
\end{quote}

\begin{center}\rule{0.5\linewidth}{0.5pt}\end{center}

\subsection{等待 GPU
(8月中旬)}\label{ux7b49ux5f85-gpu-8ux6708ux4e2dux65ec}

GPU 可用后才能执行。预计 8 月中旬出差回来、显卡到位。

\begin{itemize}
\tightlist
\item[$\square$]
  \textbf{CIFAR/MedMNIST GPU 重跑} --- ResNet-18 + 50 epochs，验证 CPU
  3-epoch 结果是否保持

  \begin{itemize}
  \tightlist
  \item
    Reason: CPU 3-epoch 太小，不足以说服审稿人
  \end{itemize}
\item[$\square$]
  \textbf{SCX-Noise 论文 Figure} --- 从重跑结果出 publication-quality
  图（ROC/PR/confusion matrix）
\item[$\square$]
  \textbf{与 Data Shapley/LESS 对比实验} --- GPU 加速后才有可比性

  \begin{itemize}
  \tightlist
  \item
    Baseline: Random, Uncertainty, Diversity, Coreset, Data Shapley
    (pyDVL), LESS, SCX
  \end{itemize}
\end{itemize}

\subsection{等待 DFT (2-4周)}\label{ux7b49ux5f85-dft-2-4ux5468}

超算运行中，等待 GaN 540 jobs + AlN 164 jobs 收敛。

\begin{itemize}
\tightlist
\item[$\square$]
  \textbf{GaN\_v1 DFT 结果 → 训练 GaN expert} --- ACE/PACE fitting，用于
  EGP Paper 1 验证
\item[$\square$]
  \textbf{AlN\_v4 DFT 结果 → 重训 AlN expert} --- 替换旧 AlN v3，做
  retrospective governance 验证
\item[$\square$]
  \textbf{EGP Paper 1 完整实验 (M1-M6)} --- gauge merging + energy
  alignment + expert compilation
\end{itemize}

\subsection{立即可做 (0
新数据)}\label{ux7acbux5373ux53efux505a-0-ux65b0ux6570ux636e}

不需要 GPU 也不等 DFT，现在就能写。

\begin{itemize}
\tightlist
\item[$\square$]
  \textbf{EGP Paper 1 Introduction + Methods + gauge fixing theory 写作}

  \begin{itemize}
  \tightlist
  \item
    方向：Introduction 已完成草稿框架，Methods 需补细节
  \item
    Priority: HIGH --- 发表优先权
  \end{itemize}
\item[$\square$]
  \textbf{SCX-Theory arXiv 草稿} --- 定义 + 命题 + 合成实验，先 GitHub
  公开再 arXiv

  \begin{itemize}
  \tightlist
  \item
    Priority: HIGH --- 抢概念占坑
  \end{itemize}
\item[$\square$]
  \textbf{合成实验运行 + 出图} ---
  \texttt{experiments/synthetic/run\_experiment.py} 已有，跑全并出
  Figure

  \begin{itemize}
  \tightlist
  \item
    预期产出：SCX four-class classification figure + value function
    visualization
  \end{itemize}
\item[$\square$]
  \textbf{SCX-Health README 完善 + GitHub 公开}

  \begin{itemize}
  \tightlist
  \item
    需完善 README，确认 Apache 2.0 LICENSE，推送到 public GitHub
  \end{itemize}
\item[$\square$]
  \textbf{Benchmark Suite 骨架} --- 标准化
  \texttt{benchmark.run(method,\ dataset)\ →\ results} API

  \begin{itemize}
  \tightlist
  \item
    5-10 数据集、Baseline 集成：Random, Uncertainty, Diversity, Shapley,
    LESS, Coreset, SCX
  \end{itemize}
\end{itemize}

\subsection{已完成 (v0.2.0 →
v0.3.0，归档至此)}\label{ux5df2ux5b8cux6210-v0.2.0-v0.3.0ux5f52ux6863ux81f3ux6b64}

以下为 2026-06-25/26 主要完成项：

\subsubsection{代码 (v0.2.0 → v0.3.0)}\label{ux4ee3ux7801-v0.2.0-v0.3.0}

\begin{itemize}
\tightlist
\item[$\boxtimes$]
  \textbf{Compress 保真定理} ---
  \texttt{theory/propositions/04\_compression\_fidelity.md} (689行)
\item[$\boxtimes$]
  \textbf{Expert Governance Protocol 形式化} ---
  \texttt{theory/propositions/05\_expert\_governance\_protocol.md}
\item[$\boxtimes$]
  \textbf{State Discovery 鲁棒性分析} ---
  \texttt{src/scx/state/robustness.py}
  (merge/split/noise/cross-descriptor/bootstrapping)
\item[$\boxtimes$]
  \textbf{自适应阈值} --- \texttt{src/scx/valuation/adaptive.py}
  (calibrate/sensitivity/auto\_threshold)
\item[$\boxtimes$]
  \textbf{StateConditionedInfluence} ---
  \texttt{valuation/influence.py}: 状态条件影响力计算
\item[$\boxtimes$]
  \textbf{OnlineSCXFramework} --- \texttt{core/online.py}: 增量状态更新
  + 在线专家追踪
\item[$\boxtimes$]
  \textbf{Clean-room 检查通过} --- 30 源文件，零学校引用
\end{itemize}

\subsubsection{实验}\label{ux5b9eux9a8c}

\begin{itemize}
\tightlist
\item[$\boxtimes$]
  \textbf{MedMNIST 3 组实验} --- compress (+6.00\% at 30\%), noise,
  routing
\item[$\boxtimes$]
  \textbf{CIFAR 3 组实验} --- baselines, noise (F1=0.617, 2.5x
  loss-based), routing
\item[$\boxtimes$]
  \textbf{两份实验报告} --- MedMNIST + CIFAR 结果文档化
\end{itemize}

\subsubsection{文档}\label{ux6587ux6863}

\begin{itemize}
\tightlist
\item[$\boxtimes$]
  \textbf{GPT 9011 行讨论整理} --- 4 份结构化文档
\item[$\boxtimes$]
  \textbf{5 篇论文谱系确立} --- 8 个文件统一更新
\item[$\boxtimes$]
  \textbf{IP 保护三文件} --- \texttt{DEVELOPMENT\_LOG.md} +
  \texttt{IP\_NOTE.md} + \texttt{CLEAN\_ROOM\_CHECK.md}
\item[$\boxtimes$]
  \textbf{SCX-Health 开源仓库} --- \texttt{scx-health/} Apache 2.0
\item[$\boxtimes$]
  \textbf{SCX\_思想扩展\_综合方案.md} --- 8 方向 + 3 阶段 roadmap
\end{itemize}

\subsubsection{项目整理}\label{ux9879ux76eeux6574ux7406}

\begin{itemize}
\tightlist
\item[$\boxtimes$]
  \textbf{4 个过期文件归档} → \texttt{CodexKnowledge/archive/}
\item[$\boxtimes$]
  \textbf{README 重写} --- 版本、结构、核心公式、关键文件索引
\item[$\boxtimes$]
  \textbf{文件夹整理} --- \texttt{paper/} 统一存放，\texttt{images/}
  集中管理
\item[$\boxtimes$]
  \textbf{EGP GaN\_v1 (540 jobs) + AlN\_v4 (164 jobs) 提交超算}
\end{itemize}

\begin{center}\rule{0.5\linewidth}{0.5pt}\end{center}

\subsection{未来方向 (v0.4.0+)}\label{ux672aux6765ux65b9ux5411-v0.4.0}

这些是旧 \texttt{SCX\_TODO.md} 和 \texttt{SCX\_思想扩展\_综合方案.md}
中未完成的后续项，暂时挂起。

\begin{itemize}
\tightlist
\item[$\square$]
  稳定性定理: state partition 小扰动 → 分类结果变化 bounded
\item[$\square$]
  PDF/PDE toy demo（悬臂梁 + Poisson PDE）
\item[$\square$]
  SCX-Semi toy demo（简化光刻模型 + OPC 路由）
\item[$\square$]
  Benchmark Suite 标准化 API 完善
\item[$\square$]
  SCX-Core 抽象接口设计（\texttt{SCXStateEncoder},
  \texttt{SCXVisionEncoder}, \texttt{SCXFEMEncoder}）
\item[$\square$]
  LLM-SCX 最小验证（LoRA 微调 0.5B 模型，20\% SCX selection vs
  baselines）
\end{itemize}

\begin{center}\rule{0.5\linewidth}{0.5pt}\end{center}

\subsection{索引}\label{ux7d22ux5f15}

\begin{longtable}[]{@{}ll@{}}
\toprule\noalign{}
主题 & 文件 \\
\midrule\noalign{}
\endhead
\bottomrule\noalign{}
\endlastfoot
开发 TODO (本文件) & \texttt{SCX\_TODO\_2026-06-26.md} \\
综合 TODO (全部方向) & \texttt{SCX\_终极TODO.md} \\
旧开发 TODO (已替代) & \texttt{SCX\_TODO.md} \\
思想扩展 (roadmap) & \texttt{SCX\_思想扩展\_综合方案.md} \\
AI 入口 & \texttt{START\_HERE\_CODEX.md} \\
\end{longtable}

\end{document}
